\documentclass{beamer}
\usetheme{Boadilla}
\usepackage[main=russian,english]{babel}

\title{Лекция 1. Введение}
\subtitle{Основы интеллектуального анализа данных}
\author{Полузёров Т. Д.}
\institute{БГУ ФПМИ}
\date{}


\begin{document}
	\begin{frame}
		\titlepage
	\end{frame}
	
	\begin{frame}
		\frametitle{Структура курса}
		\tableofcontents
	\end{frame}
		
	\section{}
	\begin{frame}
		\frametitle{Строгая зависимость}
		В мире некоторые сущности "строго" связаны. Такие зависимости можно выразить формулой
		
		\begin{itemize}
		\item Как перевести часы в минуты?
		
		
		Эта прямая функциональная зависимость.
		$f(x) = 60 * x$
		\item Как зависит ускорение тела от силы действующей на него?
		
		
		Явно выражается законом физики
		$F = m * a$
		\item Как вычислить пройденное расстояние, если известна скорость движения?
		
		Согласно формуле
		$S(T) = \int_{0}^{T}v(t) dt$
		\end{itemize}
	\end{frame}

	\begin{frame}
		\frametitle{Все ли можно выразить формулой?}
		% подогнать примеры под пункты "что может мешать"
			\textbf{А можно ли вывести формулу для нахождения курса доллара на завтра?}
			На курс влияют очень многие факторы, всех и не учесть...
			
			
			\textbf{Вычислить температуру воздуха в зависимости от других погодных условий.}
			Пусть мы наблюдаем эти признаки, но как восстановить зависимость?
			
			\textbf{Рассчитать дневное число посетителей кофейни в зависимости от погоды, дня недели?}
			Вряд ли можно выделить строгие зависимости, описать поведение единой формулой...
			
			\textbf{Узнать какая из команд выиграет футбольный матч?}
			
			\vspace{15pt}
			Что может мешать это сделать:
			\begin{itemize}
				\item Непонятно от каких факторов зависит
				\item Факторов слишком много
				\item Очень сложный вид зависимости
				\item В данных есть "неопределенность", случайность
			\end{itemize}
	\end{frame}
	
	\begin{frame}
		\frametitle{}
		Анализ данных и Машинное обучение помогает в решении этих задач:
		\begin{itemize} 
			\item Восстановление сложных зависимостей
			
			Как оценки в университете влияют на будущую ЗП?
			\item Прогнозирование 
			
			Какой курс доллара ожидается через год?
			\item Автоматизация процессов
			
			Распознавание подозрительных транзакций по банковской карте
			\item Принятие оптимальных решений 
						
			Принятие решения об выдаче кредита клиенту	
		\end{itemize}
	\end{frame}

	\begin{frame}
		\frametitle{Примеры классических задач}
			\begin{itemize}
				\item Регрессия
				\begin{itemize}
					\item Оценка стоимости квартиры по описанию
					\item Сколько товара закупить для магазина
					\item Сколько пользователей откликнется на рекламу
				\end{itemize}
			
				\item Классификация
				\begin{itemize}
					\item Выдавать ли кредит клиенту
					\item Определение наличия опухоли по рентгену
					\item К какому жанру отнести песню
				\end{itemize}
			
				\item Кластеризация
				\begin{itemize}
					\item Найти группы покупателей схожих по поведению
					\item Обнаружение "выбросов", необычных и редких явлений в данных
					\item Сжатие данных 
				\end{itemize}
			\end{itemize}
			
	\end{frame}

	\begin{frame}
		\frametitle{Постановка задачи}
		$X$ - множество объектов (objects)
		
		$Y$ - множество ответов (target)
		
		$y: X \to Y$ -неизвестная зависимость (target function)
		
		\vspace{15}
		
		\textbf{Дано}:
		
		$\{x_1, ..., x_\ell\} \subset X$ - Обучающая выборка (samples)
		$y_i = y(x_i), i=1, ..., \ell$ - известные ответы
		
		\textbf{Необходимо}:
		
		Найти алгоритм (решающую функцию, модель) $a: X \to Y$ приближающую $y$ на всем множестве $X$		
	\end{frame}

	\begin{frame}
		\frametitle{Признаковое описание}
		
		$f_j: X \to D_j, j=1, ..., n$ - признаки объекта (features)
		
		Типы признаков:
		\begin{itemize}
			\item $D_j = \{0, 1\}$ - бинарный признак $f_j$
			\item $|D_j| < \infty$ - категориальный признак $f_j$
			\item $|D_j| < \infty$, причем порядок важен - порядковый признак $f_j$
			\item $D_j \subseteq \mathbb{R}$ - количественный признак $f_j$
		\end{itemize}
	
		Вектор $(f_1, ..., f_n(x))$ - признаковое описание объекта $x$.
		
		Матрица "объекты-признаки":
		
		 $$
		 F = (f(x_{ij}))_{n \times l} = 
		 \begin{pmatrix}
		 	f_{11} & ... & f_{1\ell} \\
		 	... & ... & ... \\
		 	f_{n1} & ... & f_{n\ell} \\		 	
		 \end{pmatrix}
		 $$
	\end{frame}
		
	\begin{frame}
		\frametitle{Типы задач}
		
		Задача классификации:
		\begin{itemize}
			\item $Y = \{-1, 1\}$ - бинарная классификация
			\item $Y = \{1, ... M\}$ - на $M$ непересекающихся классов (multiclass)
			\item $Y = \{0, 1\}^{M}$ - на $M$ классов, которые могут пересекаться (multilabel)
		\end{itemize}
		
		Восстановления регрессии:
		\begin{itemize}
			\item $Y = \mathbb{R}$ или $Y = \mathbb{R}^{m}$
		\end{itemize}
	\end{frame}

	\begin{frame}
		Пару примеров задач
	\end{frame}


	\begin{frame}
		\frametitle{Модель}
		
		Модель - параметрическое семейство функций
		$$
		A = \{g(x, \theta) | \theta \in \Theta\}
		$$
		
		где $g: X \times \Theta \to Y)$ - фиксированная функция,
		$\Theta$ - множество допустимых значений $\theta$
		
		\vspace{15}
		
		\textbf{Пример:}
		
			$g(x, \theta) = \sum_{j=1}^{n} \theta_j f_j(x)$
			- линейная регрессия, $Y = \mathbb{R}$
			
			\vspace{15}
			$g(x, \theta) = sign(\sum_{j=1}^{n} \theta_j f_j(x))$
			- классификация, $Y=\{-1, +1\}\]$
	\end{frame}

	\begin{frame}
		\frametitle{Метод обучения}
		
		Метод обучения (learning algorithm) - это отображение вида
		$$
		\mu : (X \times Y)^{\ell} \to A
		$$
		
		которое произвольной выборке $X^{\ell} = (x_i, y_i)_{i=1}^{\ell}$ ставит в соответствие 
		некоторый алгоритм $a \in A$
	\end{frame}
		
	\begin{frame}
		\frametitle{Этап обучения и теста}
		\textbf{Этап обучения (train):}
		
		Метод обучения $\mu$ строит по выборке $X^{\ell} = (x_i, y_i)_{i=1}^{\ell}$ 
		алгоритм $a = \mu(X^{\ell})$
		\skip
		
		$$
		\begin{pmatrix}
			f_1(x_1) & ... & f_n(x_1) \\
			... & ... & ... \\
			f_1(x_{\ell}) & ... & f_n(x_{\ell}) \\
		\end{pmatrix}
		\xrightarrow{\text{y}}
		\begin{pmatrix}
			y_1 \\
			... \\
			y_{\ell}
		\end{pmatrix}
	% чето не работает. не видит конец формулы
		\xrightarrow{\text{\mu}}	
		a
		$$
		\textbf{Этап применения (test):}
	\end{frame}

	\begin{frame}
		\frametitle{Функционалы качества}
		
		$\mathcal{L}(a, x)$ - функция потерь (loss function) - величина ошибки алгоритма $a \in A$ на объекте $x \in X$
		
		\textbf{Функции потерь для задач классификации:}
		\begin{itemize}
			\item $\mathcal{L}(a, x) = [a(x) \not= y(x)]$ - индикатор ошибки
		\end{itemize}
		
		\textbf{Функции потерь для задач регрессии:}
		\begin{itemize}
			\item $\mathcal{L}(a, x) = |a(x) - y(x)|$ - абслолютное значение ошибки
			\item $\mathcal{L}(a, x) = (a(x) - y(x))^2$ - квадрат ошибки
		\end{itemize}
		
		Эмпирический риск - функционал качества алгоритма $a$ на $X^{\ell}$
		$$
		Q(a, X^{\ell}) = \frac{1}{\ell} \sum_{i=1}^{\ell} \mathcal{L}(a, x_i)
		$$
	\end{frame}

	\begin{frame}
		\frametitle{}
		Метод минимизации эмпирического риска:
		$$
		\mu(X^{\ell}) = arg \min_{a \in A} Q(a, X^{\ell})
		$$
		
		\textbf{Пример:} метод наименьших квадратов $Y = \mathbb{R}, \mathcal{L} - квадратична$
		$$
		\mu(X^{\ell}) = arg \min_{a \in A} \sum_{i=1}^{\ell} (g(x_i, \theta) - y_i)^{2}
		$$
		
		\textbf{Проблема обобщающей способности}
		\begin{itemize}
			\item найдем ли мы "закон природы" или переобучимся,
			то есть подгоним функцию $g(x_i, \theta) под заданные значения?$
			\item будет ли $a = \mu(X^{\ell})$ приближать функцию $y$ на всем $X$?
			\item будет ли $Q(a, X^{\ell})$ мало на новых данных - контрольной выборке $X^{k} = (x_{i}^{\prime}, y_{i}^{\prime})_{i=1}^{k}, y_{i}^{\prime} = y(x_i)$
		\end{itemize}
	\end{frame}

	% пример с переобучением пару слайдов
	
	\begin{frame}
		content...
	\end{frame}

	\begin{frame}
		\frametitle{Переобучение - одна из проблем машинного обучения}
		\begin{itemize}
			\item \textbf{Из-за чего возникает переобучение?}
			
			избыточная сложность пространства параметров $\Theta$, лишние степени свободы в модели $g(x, \theta)$ "тратятся" на чрезмерно точную подгонку под обучающую выборку. - переобучение есть всегда, когда есть оптимизация параметров по конечной (заведомо не полной) выборке.
			
			\item \textbf{Как обнаружить переобучение?}
			
			- эмпирически, путем разбиения выборки на train и test
			
			\item \textbf{Избавиться нельзя. Как минимизировать?}
			
			- минимизировать HoldOut, LOO или CV
			- наклыдвать ограничения на $\theta$ - регуляризация
			- минимизировать одну из теоретических оценок
		\end{itemize}
	\end{frame}

	\begin{frame}
		\frametitle{Эмпирические оценки обобщающей способности}
		\begin{itemize}
			\item Эмпирический риск на тестовых данных (hold-out)
			$$
			HO(\mu, X^{\ell}, X^{k}) = Q(\mu(X^{\ell}, X^{k})) \to min
			$$
			
			\item Скользящий контроль (leave-one-out) $L = \ell + 1$:
			$$
			LOO(\mu, X^{L}) = \frac{1}{L} \sum_{i=1}^{L} \mathcal{L}(\mu(X^{L} \setminus \{x_i\}), x_i) \to min
			$$
			
			\item Кросс-проверка (cross-validation), $L = \ell + k, X^{L} = X_{n}^{\ell} \sqcup X_{n}^{k}$:
			$$
			CV(\mu, X^{L}) = \frac{1}{|N|} \sum_{n \in N} Q(\mu(X_{n}^{\ell}), X_{n}^{k} ) \to min
			$$
		\end{itemize}
	\end{frame}
	
		
\end{document}